\chapter{Introduction}
\label{ch:introduction}

\epigraph{%
But how do we know when irrational exuberance has unduly escalated
asset values, which then become subject to unexpected and prolonged
contractions \dots\ And how do we factor that assessment into
monetary policy?%
}{---\,\citet{greenspan1996}}


% =============================================================================
\section{Motivation}
\label{sec:motivation}
% =============================================================================

% -----------------------------------------------------------------------------
\subsection{The recent boom in AI stocks and concentration risk}
\label{subsec:ai-stock-boom}

% beat 1: phenomenon - AI rally and concentration

On 30~November 2022, OpenAI released ChatGPT. The product
became, within weeks, a visible inflection point in market
narratives around generative AI, and was followed by a marked
revaluation of firms perceived to be central to the emerging
AI economy. At the centre of that revaluation are seven
large-capitalisation technology firms widely treated as
builders, deployers, or enablers of AI infrastructure---Apple,
Amazon, Alphabet, Meta Platforms, Microsoft, Nvidia, and
Tesla---collectively known as the Magnificent~Seven.

\Cref{fig:mag7-sp500} plots the cumulative return of an
equal-weighted Magnificent~Seven portfolio against the
S\&P~500 and the S\&P~500 with the Magnificent~Seven removed
from January~2021 through April~2026. Before late~2022, the
three portfolios moved broadly together, including through the
2022 drawdown. The divergence begins in late~2022 and
accelerates through 2023 and 2024: by early~2026 the
Magnificent~Seven had gained roughly 200\,\% relative to the
start of 2021, against around 80\,\% for the S\&P~500 as a
whole and roughly 50\,\% for the remaining 493 constituents.
The same firms have also come to occupy a much larger share of
the index itself: their combined weight in the S\&P~500 rose
from roughly 12\,\% at the start of 2015 to more than 34\,\%
by end-2025, so that seven firms now account for more than a
third of total market capitalisation \citep{motleyFoolNasdaq2026}.
Analysis by M\&G Investments attributes roughly 70\,\% of the
approximately \$24\,trillion in market value added to the
S\&P~500 since ChatGPT's launch to companies linked to AI
infrastructure and semiconductors, with around 65\,\% of the
total index rise coming from just ten stocks
\citep{vine2025}. Vontobel Asset Management reports that
42~AI-related stocks accounted for more than half of total
market returns over the same window \citep{souccar2026}.

\begin{figure}[H]
    \centering
    \includegraphics[width=\MainFigureWidth]{fig_mag7_sp500}
    \caption{Cumulative returns: Magnificent~Seven, S\&P~500, and
    S\&P~500 excluding the Magnificent~Seven, January~2021 to
    April~2026}
    \label{fig:mag7-sp500}
    \begin{minipage}{0.90\textwidth}
    \vspace{0.5em}
    \footnotesize
    \textit{Notes:} The figure plots cumulative returns for an
    equal-weighted Magnificent~Seven portfolio, the S\&P~500 index,
    and the S\&P~500 excluding the Magnificent~Seven. Sources:
    CRSP via WRDS, Investing.com, and FRED.
    \end{minipage}
\end{figure}

Concentration of this magnitude is unusual in historical
context, and it has narrowed the breadth of the rally
correspondingly: in 2023 and 2024, fewer than 28\,\% of
S\&P~500 constituents outperformed the index, the narrowest
breadth since at least~1995 \citep{firstTrust2025}. The
structural implication is that the performance of the
S\&P~500 has become unusually dependent on the valuation of a
small group of firms tied to the same technological narrative.
For investors holding passive exposure to the index through
index funds and pension products, this is a feature of the
allocation, not a choice. It also introduces substantial
downside risk: during the 2022 sell-off, the Magnificent~Seven
declined by approximately 41\,\%, roughly double the S\&P~500's
20\,\% decline. The same mechanism that amplifies gains during
a rally can magnify losses when sentiment turns, particularly
when the rally is closely tied to a single investment theme.

Whether the response is proportionate remains an open question.
Behind the surge in equity valuations lies an equally striking
surge in corporate capital expenditure.


% -----------------------------------------------------------------------------
\subsection{Capital expenditure and the valuation question}
\label{subsec:capex-valuation}
% -----------------------------------------------------------------------------
% beat 2: real-economy counterweight, valuation tension, dot-com comparison, handoff to RQs

The strong rise in AI-related stock prices raises a question
already implicit in the preceding subsection: how much of
this is a story about market enthusiasm, and how much reflects a
real change in what firms are doing? One way to bring evidence to
bear is to look at how firms are spending their money. If the
companies at the centre of the rally believe that artificial
intelligence will change how business is done, that belief should
show up in investment decisions as well as in equity prices.

\Cref{fig:ai-capex} shows that it does. Capital expenditure among
the largest U.S.\ technology firms has risen sharply, reflecting
large commitments to data centres, semiconductors, and the
infrastructure needed to train and deploy AI systems. The spending
is concentrated: Microsoft, Amazon, Alphabet, and Meta account for
the bulk of AI-related capex, and their announced plans continue
to be revised upward, with 2026 projections roughly 27\,\% above
those made a quarter earlier \citep{paydenRygel2026}. A
\citet{mckinsey2025} report projects \$6.7\,trillion in global
data-centre buildout through 2030, equivalent to more than
1\,\% of global GDP annually, with about half allocated to
servers and chips. Nvidia's sales rose from \$24\,billion in
2021 to \$187\,billion in 2025, driven almost entirely by demand
from this same cluster of firms \citep{mercer2025}. The
infrastructure investment is therefore not only large in absolute
terms but also concentrated in the same firms whose share prices
dominate the S\&P~500.

\begin{figure}[H]
    \centering
    \includegraphics[width=\MainFigureWidth]{fig_ai_capex}
    \caption{Capital expenditure among major technology firms, 2015--2027E}
    \label{fig:ai-capex}
    \begin{minipage}{0.90\textwidth}
    \vspace{0.5em}
    \footnotesize
    \textit{Notes:} The figure reports capital expenditure among
    major technology firms. Sources: SEC filings, FactSet consensus
    estimates, and \citet{golan2025}.
    \end{minipage}
\end{figure}

The scale of the spending does not on its own resolve the
valuation question; if anything, it sharpens it. Unlike in the
late-1990s technology cycle, the current buildout has been largely
self-financed: U.S.\ public technology companies have funded their
AI capital expenditure almost entirely from free cash flow, with
the aggregate ratio of free cash flow to capital spending across
the top 3{,}000 U.S.\ stocks remaining below one as of late~2025,
against a dot-com peak of nearly four \citep{ojjehFronczke2026}.
Within the Magnificent~Seven, the ratio of operating expenditure
to capital expenditure has fallen from a long-run range of
2.5--4$\times$ to roughly 1.5$\times$ since mid-2024, and return
on invested capital at Meta and Alphabet has remained above
30\,\% across the period \citep{ojjehFronczke2026}. Self-financing
reduces the risk of a credit-driven unwind, but it does not
eliminate the risk that the revenue these investments are
expected to generate may not materialise at the scale or pace
markets currently price in. OpenAI projects revenues rising from
\$13\,billion in 2025 to \$200\,billion in 2030 against cumulative
cash burn of roughly \$119\,billion before profitability
\citep{mercer2025}, and competition from efficient open-source
models has begun to put downward pressure on pricing.

The natural comparison is to the late-1990s technology cycle,
where rapid price appreciation in the leading firms ran ahead of
realised earnings and ultimately reversed.
\Cref{fig:price-earnings-comparison} shows the comparison
directly: in the dot-com era, the share prices of the leading
technology firms rose far faster than their forward-earnings
outlook, while in the current AI episode the rise in share
prices has been more closely accompanied by stronger earnings
growth. This does not imply that current valuations are justified
by fundamentals---earnings projections themselves may embed
optimistic assumptions about adoption, monetisation, and future
market power, and the same projections are part of what supports
the current pricing. But the gap between prices and the
contemporaneous earnings outlook is narrower than it was in the
final phase of the dot-com boom, and the historical parallel is
therefore weaker than a price-only comparison would suggest.

\begin{figure}[H]
    \centering
    \begin{subfigure}[t]{0.49\textwidth}
        \centering
        \includegraphics[width=\linewidth]{fig_pe_dotcom}
        \caption{Dot-com era}
        \label{fig:pe-dotcom}
    \end{subfigure}
    \hfill
    \begin{subfigure}[t]{0.49\textwidth}
        \centering
        \includegraphics[width=\linewidth]{fig_pe_ai}
        \caption{AI era}
        \label{fig:pe-ai}
    \end{subfigure}
    \caption{Stock prices and forward earnings in the dot-com era
    and the current AI era}
    \label{fig:price-earnings-comparison}
    \begin{minipage}{0.95\textwidth}
    \vspace{0.5em}
    \footnotesize
    \textit{Notes:} The figure compares the evolution of share
    prices and forward earnings for leading technology firms in
    the dot-com era (panel a) and the current AI era (panel b).
    \end{minipage}
\end{figure}

What remains is a valuation question with two empirical sides.
The first concerns the price series themselves: does the
behaviour of AI-exposed equities, looked at across the
cross-section of exposure, show statistical signatures that go
beyond what fundamentals-driven pricing would imply? The second
concerns how those same prices respond to news that should
matter most for firms with cash flows far in the future: are
AI-exposed firms differentially sensitive to the marginal
monetary-policy information that markets read as relevant for
long-duration valuations? These are the two questions the rest
of the thesis takes up.


% =============================================================================
\section{Research questions and findings}
\label{sec:intro-rq}

% beat 3a+4a: PSY question + method
% Note: this section's two questions are introduced; the soft-form
% handoff at the end of subsec:capex-valuation already frames them,
% so we go straight into the first question here (no separate intro
% paragraph) to avoid restating the pair a third time.
The first question asks whether the price behaviour of
AI-exposed equities shows statistical signatures of exuberance
once one looks beyond the Magnificent~Seven, across the
cross-section of disclosed AI exposure. The thesis answers this
by applying the recursive right-tailed unit-root procedure of
\citet{phillipsShiYu2015a}---the PSY test, after the authors,
whose test statistic is the generalised sup-ADF (GSADF)
statistic---to value- and equal-weighted total-return
portfolios of S\&P~500 firms partitioned by their pre-2023 10-K
AI keyword score, together with within-equity-universe and
broad-market relative ratios and three AI-concentration baskets,
following the price-to-index ratio specification of
\citet{baselePhillipsShi2025}. The classification is fixed on
pre-2023 disclosures, before the rally it is used to study, so
the partition is not constructed with hindsight. Rejection of
the unit-root null is read as evidence of mildly explosive price
dynamics, not as evidence of mispricing.

% beat 3b+4b: LP question + method
The second question asks whether AI-exposed firms are
differentially sensitive to the marginal monetary-policy news
that markets read as relevant for long-duration cash flows. The
thesis answers this by estimating panel local projections
\citep{jorda2005} of firm-level cumulative log returns on the
\citet{swanson2021} forward-guidance/path-factor surprise,
interacted with AI-exposure dummies, at horizons $h = 0$ to
$60$ trading days after each scheduled FOMC announcement.
Firm fixed effects and GICS-sector\,$\times$\,event fixed
effects absorb every component of the response that is common
to firms in the same sector on the same FOMC date, so the
coefficients of interest identify a within-sector, within-event
differential response of High~AI and Moderate~AI firms relative
to No~AI firms, rather than a level effect.


% -----------------------------------------------------------------------------
% beat 5: findings preview - the two paragraphs already drafted
% -----------------------------------------------------------------------------

% beat 5a: PSY findings preview
The first analysis tests for explosive price dynamics in
value- and equal-weighted portfolios of S\&P~500 firms sorted
by pre-2023 10-K AI exposure, using the recursive right-tailed
unit-root procedure of \citet{phillipsShiYu2015a} on the
price-to-index ratio specification of \citet{baselePhillipsShi2025}.
Over the monthly sample 2010-01 to 2026-04, the High AI / No AI
ratio rejects the unit-root-with-drift null at the 99\,\% level
under both weightings (GSADF $3.15$ value-weighted, $3.19$
equal-weighted). The strongest concentration signal lies in the
Mag-7 / No AI ratio (GSADF $5.20$ value-weighted and $6.72$
equal-weighted, both at the 99\,\% level), with the
Mag-7 / High AI ratio also rejecting at 99\,\% under both
weightings. The Moderate AI series rejects only at the 95\,\%
level (GSADF $2.35$ value-weighted, $2.24$ equal-weighted) and
yields a single short 2021 episode under each weighting,
placing it closer to the No AI benchmark than to High AI.
Date-stamping qualifies the headline: the broad High AI / No AI
ratio is flagged mainly in pre-ChatGPT episodes (2018; 2020--21),
while the post-2022 signal appears most clearly in the
value-weighted High-AI absolute index and in the Mag-7
concentration baskets relative to both denominators. The
evidence is consistent with bubble-like dynamics around
AI-exposed and AI-associated firms, but it is not a clean
post-2022 AI-bubble verdict: the High AI cohort overlaps with
a longer-running pattern of technology-sector outperformance,
and rejection of the unit-root null is evidence of mildly
explosive trajectories, not of mispricing per se.

% beat 5b: LP findings preview
The second analysis estimates the differential equity-return
response of AI-exposed firms to forward-guidance monetary-policy
shocks. Panel local projections \citep{jorda2005} of cumulative
log returns on the \citet{swanson2021} path-factor surprise
interacted with AI-exposure dummies, with firm and
sector\,$\times$\,event fixed effects and two-way clustered
standard errors, deliver a horizon profile that is precisely
zero on impact and significantly negative at long horizons.
For High AI firms relative to No AI firms in the same GICS
sector on the same FOMC date, a $+25$ basis-point hawkish path
surprise is associated with a cumulative return $6.49$
percentage points lower at $h = 40$ trading days
(SE $2.11$, 95\,\% CI $[-10.62, -2.35]$) and $7.60$ percentage
points lower at $h = 60$ (SE $2.53$, 95\,\% CI $[-12.56, -2.63]$).
At shorter horizons ($h \leq 20$) the differential is small in
absolute value and the 95\,\% confidence interval contains zero
at every key horizon. The Moderate AI differential is small at
every key horizon and not statistically distinguishable from
zero, so the pattern is tail-versus-rest rather than a graded
dose-response in measured AI exposure. The differential
survives alternative sample windows, alternative AI
classifications, abnormal-return outcomes, and a partition of
the path shock by the sign of contemporaneous S\&P futures
reaction \citep{bauerSwanson2023}; a leverage-amplification
check rules out financial leverage as the source of the
heterogeneity. The estimand is a within-sector, within-event
differential rather than a level effect, and the design does
not separate discount-rate, cash-flow, and risk-premium
channels.


% =============================================================================
\section{Contribution}
\label{sec:intro-contribution}

% beat 6a: contribution

The thesis contributes to two literatures and one combination of them.

First, it adds AI exposure as a firm-level heterogeneity dimension
in the literature on monetary transmission to equity prices.
That literature has so far focused on financial and structural
firm characteristics, with leverage and debt structure receiving
the most attention \citep{ottonelloWinberry2020,
gurkaynakKarasoyCanLee2022, cloyneEtAl2023, arinEtAl2025}; AI
exposure is not currently among them. The differential identified
in \cref{ch:results} survives a leverage-amplification check
(\cref{sec:res-heterogeneity}), so it is not a relabelled
financial-leverage differential and stands as a candidate
characteristic for inclusion in this literature alongside the
established financial dimensions.

Second, it extends the bubble-detection evidence on the recent AI
boom along a cross-sectional and relative-pricing dimension.
\citet{baselePhillipsShi2025} apply the PSY procedure to the
Nasdaq aggregate and to the seven individual Magnificent~Seven
stocks; the thesis applies the same procedure and bootstrap
framework to value- and equal-weighted portfolios of S\&P~500
firms partitioned by their pre-2023 disclosure-based AI exposure,
together with within-equity-universe and broad-market ratios and
three AI-concentration baskets. This locates the source of
explosive dynamics within the cross-section of AI engagement,
rather than reading them off a single index or a fixed mega-cap
basket.

Third, the most distinctive contribution comes from holding the
same AI-exposure partition fixed across two non-nested empirical
exercises on a common sample window. The local-projection
exercise alone could be read as evidence that AI-exposed firms
inherit the discount-rate sensitivity of growth firms documented
in the prior asset-pricing literature; the PSY exercise alone
could be read as a cross-sectional extension of an
already-documented Magnificent~Seven phenomenon. The joint
finding---that the same partition along which the
shock-conditional differential is strongest is also the partition
along which explosive price-level dynamics are most
pronounced---is the substantive content of the thesis. It is
reported as a pattern across two designs, not as an identified
mechanism.


% =============================================================================
\section{Roadmap}
\label{sec:intro-roadmap}

The remainder of the thesis is organised as follows.
\Cref{ch:literature} reviews the literatures the empirical
strategy builds on---high-frequency identification and
forward-guidance shocks; the equity-valuation channels of a
present-value decomposition; firm-level heterogeneity in
monetary transmission; the AI-investment and intangibles
literature; and the recursive right-tailed unit-root testing
framework---and locates the thesis's contribution within them.
\Cref{ch:data} describes the S\&P~500 firm universe, the
pre-2023 10-K AI-exposure classification shared by both
analyses, and the two analysis-specific data builds.
\Cref{ch:empirical_strategi} sets out the panel local-projection
design and the PSY recursive testing design, including the
identification limits of each. \Cref{ch:results} reports the
main empirical results. \Cref{ch:robustness} reports the
specification, sample, and inference robustness exercises for
both analyses. \Cref{ch:conclusion} concludes.